On injecte à un patient un médicament. Après l'injection, on mesure la concentration de ce médicament dans le sang, en grammes par litre (g/L), pendant 13 heures.

Cette concentration est modélisée par la fonction $f$ définie sur l'intervalle $[0 ; 13]$ par :

$f(t)=\left(t^{2}+t+1\right) \e^{-t}$ \quad où $t$ représente le nombre d'heures écoulées depuis l'injection.

\begin{enumerate}
	\item Déterminer $f(0)$ et interpréter ce résultat dans le contexte de l'exercice.
	\item\begin{enumerate}
		\item Justifier que pour tout réel $t$ dans $[0 ; 13]$, on a : \quad $f'(t)=t(-\mathrm{t}+1) \e^{-t}$.

		\item Étudier le signe de $t(-\mathrm{t}+1)$ sur $[0 ; 13]$.

		\item En déduire le tableau de variations de $f$ sur $[0 ; 13]$.

		\emph{Il n'est pas demandé de calculer les images.}
		\end{enumerate}

	\item On estime que le médicament n'est plus actif lorsque sa concentration est strictement inférieure à 0,1~g/L. Il convient alors de faire une nouvelle injection du médicament. En utilisant le tableau de valeurs approchées ci-dessous, déterminer le nombre d'heures maximal $t_{0}$ ($t_{0}$ est un entier naturel) pendant lequel le médicament reste actif.
\end{enumerate}

\begin{tblr}{colspec={*{12}{X[c]}},hlines, vlines, column{1}={gray8}}
	$t$ & 3 & 4 & 5 & 6 & 7 & 8 & 9 & 10 & 11 & 12 & 13 \\
	$f(t)$ & 0,647 & 0,385 & 0,209 & 0,107 & 0,052 & 0,025 & 0,011 & 0,005 & 0,002 & 0,001 & 0 \\
\end{tblr}


\needspace{10\baselineskip}
